% Generated draft: requires mathematical and source review.
\documentclass[11pt]{article}
\usepackage[margin=25mm]{geometry}
\usepackage{fontspec}
\setmainfont{Latin Modern Roman}
\usepackage{amsmath,amssymb,mathtools,unicode-math}
\setmathfont{Latin Modern Math}
\usepackage{enumitem,needspace,xurl,hyperref}
\setlength{\parindent}{0pt}
\setlength{\parskip}{5pt}
\setlength{\emergencystretch}{3em}
\newcommand{\N}{\mathbb{N}}
\newcommand{\Z}{\mathbb{Z}}
\newcommand{\Q}{\mathbb{Q}}
\newcommand{\R}{\mathbb{R}}
\newcommand{\C}{\mathbb{C}}
\newcommand{\F}{\mathbb{F}}
\newcommand{\A}{\mathbb{A}}
\newcommand{\PP}{\mathbb{P}}
\newcommand{\E}{\mathbb{E}}
\newcommand{\eps}{\varepsilon}
\newcommand{\dd}{\,\mathrm{d}}
\begin{document}
\section*{UnsolvedMath ID: 30003815}
\textbf{Source problem code:} catalogue entry
\subsection{Definitions and mathematical statement}
For a positive integer \(m\), an ordinary alternating-sign matrix of order \(m\) is a matrix \(M=(m_{ij})\in\{-1,0,1\}^{m\times m}\) such that every row and column sums to \(1\), and the nonzero entries in each row and column alternate in sign. Denote the set of these matrices by \(\operatorname{ASM}_m\).

The problem asks for a natural definition of a family \(\mathcal A_{n,k}\) of objects called \((n,k)\)-alternating-sign matrices. Write \(\mathcal C_{n,k}\) for the intended class of \((n,k)\)-column-strict shifted plane partitions mentioned in the source. The requested properties are
\[
\lvert\mathcal A_{n,k}\rvert=\lvert\mathcal C_{n,k}\rvert
\qquad\text{for every admissible pair }(n,k),
\]
and
\[
\mathcal A_{n-1,2}=\operatorname{ASM}_n
\quad\text{or there is a simple bijection}\quad
\mathcal A_{n-1,2}\longleftrightarrow\operatorname{ASM}_n,
\]
for every \(n\) in the intended range of this specialization.

These formulas describe the requested construction, subject to unresolved source qualifications. The supplied record does not define \(\mathcal C_{n,k}\), specify the domains of \(n\) and \(k\), or give the shape, entry bounds, and parameter restrictions governing these plane partitions. Thus \(\mathcal C_{n,k}\) is a placeholder for the intended source class, not a definition supplied here. The admissible ranges must be recovered before the displayed requirements become a fully specified statement. In particular, whether \(n=1\) is included, thereby requiring \(\mathcal A_{0,2}\), is unresolved. The requirements that the definition be natural and the bijection simple are qualitative and are not formalized by the supplied record.
\subsection{Short English statement}
Define a natural family of \((n,k)\)-alternating-sign matrices equinumerous with the intended \((n,k)\)-column-strict shifted plane partitions, whose \((n-1,2)\) specialization recovers ordinary alternating-sign matrices of order \(n\), either directly or through a simple bijection. The parameter ranges and the precise plane-partition class are missing from the supplied record.
\subsection{Sources}
\begin{itemize}
\item Local UnsolvedMath registry: permanent ID \(30003815\), source problem code \texttt{catalogue entry}, title \emph{Generalized Alternating-Sign Matrices and Shifted Plane Partitions}, in the Oberwolfach Reports Open Problems collection. No individual proposer is identified in the supplied record.
\item Supplied report citation: \emph{Enumerative Combinatorics (2019)}, Oberwolfach Reports, pp.~1381--1464, DOI \texttt{10.4171/owr/2018/23}; \url{https://doi.org/10.4171/owr/2018/23}. The background attributes the collection to mathematical workshop reports and their open-problem discussions. The literature-source list labels this same link ``original OWR report.'' No problem-specific page is supplied.
\item Literature assessment dated exactly \texttt{2026-08-21}: the supplied curation labels the problem open and reports that no general natural definition satisfying the required equality was located in the available literature. This is supplied dated curation, not independently verified source consultation or a verified assertion of current mathematical status. The supplied record also labels the statement \texttt{corrected\_verified}, explaining that notation, typography, and sentence structure were repaired while mathematical content was preserved; that label is not an independent verification performed for this draft.
\end{itemize}
\end{document}
